% Transcription — fonds Alexandre Grothendieck, shelfmark 67, batch 5 (pages 81-100).
% Established from the facsimile archives/batches/67/batch-05.pdf (a mirror of
% https://grothendieck.umontpellier.fr/67.pdf), per the skill
% transcribe-grothendieck. The batch file carries no cover sheet: PDF page N
% is page 80+N, as the Montpellier footer printed on each PDF page confirms.
% \page{} numbers that position, as batch 4 does; the archivists' pencil
% numbers bottom-left run three ahead (page 81 is pencilled 84, page 100 is
% pencilled 103).
%
% Pass: Opus 5.5 (claude-opus-5-5), 2026-09-22 — first pass, unchecked against the pages by a human.
% Revised 2026-09-23 (Opus 5.5), against the facsimile: the header no longer
% says the torsors E^!_a give surfaces « with boundary » — p. 100 writes
% « surfaces conformes sans bord », « sans » apparently crossed out and
% nothing written in its place; the « à » the first pass read there is not
% on the page, and the body of p. 100 now says so.
%
% What the batch is.
%   81      Two lines only, on a leaf he numbers « 15 »: the end of the
%           sentence batch 4 breaks off on (page 80, « il ne devrait pas
%           être bien dur de vérifier si »).
%   82, 89  Blank. Absent.
%   83-88   « Modules M_{0,4} et l'octaèdre » (his boxed title). I a set of
%           four points, J = P_{2,2}(I), V(J) the Klein four-group; the
%           plane V ⊂ S_I acting on a projective line X, its six fixed
%           points forming an octahedron S; the normaliser Γ ≃ S_4, the
%           exact sequence 1 → V → Γ → Aut(V) → 1 and its splittings as
%           pairs of antipodal faces; a Proposition (p. 85) matching
%           Klein subgroups V ⊂ Aut X, subgroups Γ ≃ S_4, and octahedral
%           6-subsets S ≃ (P^1, {0, ∞, ±1, ±i}); then the V-torsor I(V),
%           the formulas X = X(I) ∧^V t and X/V ≃ Σ(J) (p. 87), and a
%           Théorème (p. 88): the groupoid of (X, I ↪ X), I of card. 4,
%           is equivalent to that of pairs (I, ξ ∈ Σ(J(I))*).
%   90-100  « Courbes elliptiques réelles », dated « Déc 83 » top left of
%           p. 90. First (pp. 90-91) the linear algebra: a 2-regular
%           module M with an involution σ ↔ triples (P, Q, m) with
%           P = M_+, Q = M_-, m ⊂ P/2P ⊕ Q/2Q, and a Prop. that this is an
%           equivalence; p. 91 is the rough calculation for it. Then
%           (pp. 92-97) real elliptic curves as (Π, σ, c): the two
%           families E_a, E'_a (a ∈ R*+), Q(x e_+ + y e_-) = x² + a y²,
%           real structures on a complex curve as a torsor under G^+, the
%           two G^+-orbits, a(−σ) = 1/a(σ), and the case analysis n = 2,
%           4 (square: E_1, E'_1), 6 (hexagonal: E'_{√3}, E'_{1/√3}).
%           P. 95 is a leaf of drawings only (spheres cut by arcs).
%           P. 98 gives τ = √a·i for E_a, τ = ½(1 + √a·i) for E'_a, and
%           counts components of X(R): two for E_a, one for E'_a (p. 100).
%           P. 99 returns to the module classification over a principal
%           ring (conjugacy classes of involutions in GL(2, k)). P. 100
%           closes with H¹(R, E) via 0 → Π → V → E(C) → 0: Z/2 for E_a,
%           0 for E'_a, hence the only genuinely inhomogeneous real
%           elliptic curves are the torsors E^!_a; the page says they give
%           the non-orientable conformal surfaces of genus 1, « sans bord »
%           with « sans » seemingly crossed out and no word put in its place,
%           so whether he meant with or without boundary is not settled by
%           the page (the modernised reading shows it is without).
%           The margin of p. 100 carries « 9.2. », perhaps a date.
%
% Notation fixed for the batch. The symmetric group is \mathfrak{S}; the
% set of 2-element (resp. 3-, 6-element) subsets is \mathfrak{P}_{2}(S)
% etc.; blackboard letters as batch 4 has them (\mathbb{R}, \mathbb{C},
% \mathbb{Z}, \mathbb{P}); the underlined antipodism a is \underline{a},
% the underlined torsor t is \underline{t}, both in math only. His
% contracted product X(I) ∧_V t, V written under the wedge, is set
% \wedge^{V}. The letter he draws as an X barred horizontally, in Σ(J)
% and Σ(J)*, is read \Sigma: it is the quotient X/V, a line.
%
% Legibility. Pages 83-85, 87-88 and 97-100 are written fair and read
% nearly through; 86, 92, 93 and the lower half of 96 are a fast draft
% heavily struck over, whose formulas carry and whose prose often does
% not. The oblique margins of 83 and 85 are largely lost even turned.
\documentclass[11pt,a4paper]{article}
% Le préambule commun à toutes les transcriptions du fonds Grothendieck.
%
% Il tient en une idée : **l'incertitude se compose, elle ne se résout pas.**
% Une transcription qui lisse un mot douteux a détruit la seule chose pour
% laquelle on la faisait — un lecteur doit pouvoir savoir, à chaque endroit,
% ce qui était sur le feuillet et ce qui a été ajouté. Les six macros qui
% suivent sont donc l'appareil critique, et non de la décoration.
%
% Les mêmes commandes sont reconnues par `scripts/render.mjs`, qui produit la
% vue de lecture du site. Une macro ajoutée ici doit l'être aussi là-bas,
% faute de quoi le rendu échouera bruyamment — ce qui est voulu : un rendu
% silencieusement faux est pire qu'un rendu absent.

\ProvidesPackage{grothendieck}[2026/08/08 Transcriptions du fonds Grothendieck]

\usepackage{amsmath,amssymb,amsthm}
\usepackage{mathrsfs}
\usepackage{stmaryrd}
% Les diagrammes commutatifs sont partout dans ce fonds : c'est la langue
% même de la géométrie algébrique de Grothendieck, et une transcription qui
% les rendrait en prose ne transcrirait rien.
\usepackage{tikz-cd}
% Français par défaut — c'est la langue des feuillets — mais l'anglais est
% chargé aussi : la traduction et le résumé sont des documents anglais, et la
% typographie française (espaces avant les deux-points) y serait une faute.
% Ces documents basculent par \selectlanguage{english} en tête de corps.
\usepackage[english,french]{babel}
\usepackage{fontspec}
\usepackage{geometry}
\geometry{a4paper,margin=25mm,marginparwidth=28mm}
\usepackage{marginnote}
\usepackage{xcolor}
% ulem, not soul. `soul`'s \st and \ul are built on a character-by-character
% scan that XeLaTeX's font machinery breaks: the document fails with an
% undefined control sequence rather than degrading. ulem does the same two jobs
% and is Unicode-engine safe. `normalem` stops it hijacking \emph, which the
% transcriptions use for Grothendieck's own emphasis.
\usepackage[normalem]{ulem}
\usepackage{hyperref}
\hypersetup{colorlinks=true,linkcolor=black,urlcolor=black,pdfborder={0 0 0}}

% --- Métadonnées du lot -----------------------------------------------------
% Reprises telles quelles par le script de rendu, qui les affiche en tête de
% la vue de lecture. Elles fixent ce que le fichier prétend couvrir : sans
% elles, une transcription flotte sans point d'attache dans un fonds de seize
% mille feuillets.

\newcommand{\folder}[1]{\gdef\grothFolder{#1}}
\newcommand{\batch}[1]{\gdef\grothBatch{#1}}
\newcommand{\leaves}[2]{\gdef\grothFirst{#1}\gdef\grothLast{#2}}
\newcommand{\foldertitle}[1]{\gdef\grothFoldertitle{#1}}
\gdef\grothFolder{?}\gdef\grothBatch{?}\gdef\grothFirst{?}\gdef\grothLast{?}\gdef\grothFoldertitle{}

% --- Le résumé --------------------------------------------------------------
%
% Il ouvre la lecture modernisée, sous le titre « Résumé », et il est composé
% en petit. Ce n'est pas un ornement : c'est le seul passage du document qui
% ne soit pas des mathématiques, et un lecteur doit pouvoir le reconnaître
% avant de l'avoir lu — puis le sauter s'il connaît déjà le sujet, ou s'y
% arrêter s'il ne le connaît pas. Le filet qui le suit marque la reprise du
% corps.

\newenvironment{resume}
  {\small\setlength{\parskip}{0.45em}}
  {\par\medskip\hrule\medskip}

% --- Mots-clés ---------------------------------------------------------------
%
% En anglais, en clôture du résumé de la lecture modernisée : le vocabulaire
% moderne sous lequel chercher ce que les feuillets construisent. Le manifeste
% du site (`npm run manifest`) les extrait pour étiqueter la cote entière —
% cette ligne est la source unique des tags, il n'y a pas de fichier à côté.
\newcommand{\keywords}[1]{\par\smallskip\noindent
  {\footnotesize\textsc{Keywords} --- \textit{#1}}\par}

% --- L'appareil critique ----------------------------------------------------

% Le numéro de feuillet, en marge. C'est l'ancre qui permet de régler un
% désaccord sur une formule en regardant *un* feuillet plutôt qu'une chemise
% de six cents. Le site s'en sert aussi pour tourner les pages du fac-similé
% quand on fait défiler la transcription.
\newcommand{\leaf}[1]{%
  \marginnote{\footnotesize\color{gray!70}\textsf{#1}}%
  \label{leaf:#1}\ignorespaces}

% Illisible : on ne devine pas. Un mot inventé qui se lit comme les autres est
% le pire résultat possible.
\newcommand{\ill}{\textcolor{red!60!black}{[\,\ldots\,]}}

% Lecture douteuse : le mot est donné, et signalé comme incertain.
\newcommand{\uncertain}[1]{\uline{#1}}

% Ajout de l'éditeur — une lettre manquante, un mot sous-entendu.
\newcommand{\add}[1]{\textcolor{blue!50!black}{[#1]}}

% Biffé par Grothendieck lui-même. Ce qu'il a rayé fait partie du document :
% une rature dit souvent où la pensée a changé de direction.
\newcommand{\struck}[1]{\sout{#1}}

% Note du transcripteur, sur l'état matériel du feuillet.
\newcommand{\note}[1]{\par\smallskip{\small\color{gray!60!black}\textit{[#1]}}\par\smallskip}

% Note marginale de Grothendieck, distincte du corps du texte.
\newcommand{\marginal}[1]{\par\smallskip\noindent\hspace{1em}%
  {\small\color{gray!60!black}#1}\par\smallskip}

% --- Titre ------------------------------------------------------------------

\AtBeginDocument{%
  \begin{center}
    {\large\bfseries Fonds Alexandre Grothendieck --- cote n\textsuperscript{o}~\grothFolder}\\[2pt]
    {\itshape\grothFoldertitle}\\[4pt]
    {\small feuillets \grothFirst--\grothLast\ (lot \grothBatch)}\\[2pt]
    {\footnotesize\color{gray!60!black}Transcription établie d'après le fac-similé numérisé
     par l'Université de Montpellier.\\
     Les lectures incertaines sont soulignées, les illisibles notées [\,\ldots\,],
     les ajouts entre crochets.}
  \end{center}
  \vspace{1em}}

\endinput


\folder{67}
\batch{5}
\pages{81}{100}
\dating{[à partir de 1977]-1984}
\watermark{Édition de démonstration}
\foldertitle{Chirurgie des surfaces conformes : notes manuscrites (s.d., 1983-1984), lettres (1984)}

\begin{document}

\page{81}
\note{feuillet numéroté « 15 » en tête par l'auteur, qui n'y écrit que deux
lignes : elles achèvent, semble-t-il, la phrase sur laquelle s'arrête la
page 80 (lot 4), « il ne devrait pas être bien dur de vérifier si »}

est \uncertain{étale}, en calculant \uncertain{ou} interprétant
\uncertain{l'application} tangente.

\section{Modules $M_{0,4}$ et l'octaèdre}
\note{titre encadré, en tête du feuillet, de sa main}

\page{83}
$I$ ens.\ de card.\ 4 donné

$J = \mathfrak{P}_{2,2}(I)$, $V(J)$ plan vectoriel sur $\mathbb{F}_{2}$ avec
$V(J)^{*} \simeq J$,

$V(J) \subset \mathfrak{S}_{I}$ formé des \uncertain{bi-identités} et des
« bitranspositions ».

On s'intéresse aux droites projectives (sur $\mathbb{C}$ disons -- \ill{} sur
tt corps alg.\ clos \add{de car.\ $\neq 2$}) munies de $I \hookrightarrow X$.
(plus précisément $I_{\mathbb{C}} \hookrightarrow X$, ou $I \to X(\mathbb{C})$
…).
\note{« de car.\ $\neq 2$ » est ajouté en interligne}

L'opération de $V(J)$ sur $I$ se prolonge en une opération de $V(J)$ sur $X$
\[
  V(J) \hookrightarrow G_{X} = \operatorname{Aut}(X)
\]
Soit $S$ = ens.\ des pts dans $X$ qui sont fixes par un élément de
$J = V(J)^{*}$, donc $\operatorname{card} S = 6$.

\marginal{\ill{} \uncertain{ok} \ill{}}
[En fait, on trouve que les opérations \uncertain{fidèles} d'un plan vectoriel
sur une $X$, i.e.\ les plongements $V \hookrightarrow G_{X}$, sont toutes
équivalentes (modulo autom.\ intérieurs de $G_{X}$) et sont \uncertain{même}
\uncertain{avec base fixée}${}^{e_{0},\, e_{1}}$ équivalentes : l'opération
sur $\mathbb{P}^{1}$
\[
\begin{array}{lll}
  u(e_{0}) = (z \mapsto -z) = \begin{pmatrix} -1 & 0 \\ 0 & 1 \end{pmatrix}
    & & \text{pts fixes } 0, \infty \\[1ex]
  u(e_{1}) = (z \mapsto 1/z) = \begin{pmatrix} 0 & 1 \\ 1 & 0 \end{pmatrix}
    & & \text{pts fixes } \pm 1 \\[1ex]
  u(e_{0}+e_{1}) = (z \mapsto -1/z) = \begin{pmatrix} 0 & -1 \\ 1 & 0 \end{pmatrix}
    & & \text{pts fixes } \pm i
\end{array}
\]

Les \struck{quatre groupes} \uncertain{sous-groupes} \ill{} : Donc les
\ill{} de $G_{X}$ correspondent 1-1 aux parties $S$ de cardinal 6 de $X$,
isom.\ à $(\mathbb{P}_{1}, \{0, \infty, \pm 1, \pm i\})$.
\marginal{parties \uncertain{octaédrales} \uncertain{d'une droite}
\uncertain{projective}}
Le normalisateur d'un tel \struck{\ill{}} sous-groupe $V$ \add{de $G_{X}$}
\uncertain{est isomorphe} au groupe \struck{des autom.} $\mathfrak{S}_{4}$ (de
façon \uncertain{canonique}, i.e.\ à $\mathfrak{S}_{I}$ \ill{} $V$
\ill{}), plus précisément, si $V = V(J)$ avec $J = \mathfrak{P}(I)$, on
trouve \struck{des} \uncertain{un} iso de suites exactes
\note{la marge gauche porte, écrit en oblique, « parties octaédrales d'une
droite projective » (lecture douteuse) ; plus haut, à hauteur du crochet
ouvrant « [En fait », une autre note oblique de trois ou quatre mots qui ne se
lit pas}

\begin{tikzcd}[column sep=small]
  1 \arrow[r] & V \arrow[r] \arrow[d, no head, "\Vert" description] & \mathfrak{S}_{I} \arrow[r] \arrow[d] & \operatorname{Aut}(V) \arrow[r] \arrow[d, no head, "\Vert" description] & 1 \\
  1 \arrow[r] & V \arrow[r] & \Gamma \arrow[r] & \operatorname{Aut}(V) \arrow[r] & 1
\end{tikzcd}

\note{au-dessus de $\mathfrak{S}_{I}$, un mot biffé et, sous lui, un signe
d'égalité vertical}
défini modulo automorphismes intérieurs par un élément de $V$. Quand $I$ est
le torseur trivial, on trouve un tel \uncertain{iso} d'ex\ill{}

\page{84}
précisé : se donner un scindage de l'extension
\[
  (*) \qquad 1 \to V \to \Gamma \to \operatorname{Aut}(V) \to 1,
  \qquad \Gamma = \operatorname{Nor}_{G}(V), \quad
  \operatorname{Aut}(V) \simeq \mathfrak{S}_{J}
\]
\note{« $\operatorname{Nor}_{G}(V)$ » est écrit sous $\Gamma$ avec un signe
d'égalité vertical, et « $\simeq \mathfrak{S}_{J}$ » sous
$\operatorname{Aut}(V)$ ; l'indice de Nor est une petite lettre qui se lit
mal}
En termes géométriques, ces scindages \uncertain{sont associés} aux quatre
systèmes de \uncertain{paires de} \uncertain{faces} antipodiques sur
l'octaèdre. \struck{Ainsi les} L'octaèdre combinatoire (non orienté) \ill{}
l'ens.\ $S$ de \ill{} \uncertain{points}) se définit : partition de \ill{} et
l'antipodisme \struck{sur} $S$, de la façon \ill{}, et \ill{} habituelles,
on forme les parties : trois éléments \struck{qui} $S'$ telles que
$S' \cup \underline{a}(S') = S$. \struck{Pour} un tel triplet, le sous-groupe
de $\Gamma$ qui l'invarie est un scindage de l'extension $(*)$.

\ill{} le groupe $\Gamma$ \struck{est} \uncertain{est} le groupe des autom.\ de
la \uncertain{structure octaédrale} \struck{\ill{}} $S$ qui conservent
\struck{\ill{}} \uncertain{les faces} $F \subset \mathfrak{P}_{3}(S)$ \ill{}.
\uncertain{Orientation} : (définie par l'ens.\ des \ill{}
\note{« $F \subset \mathfrak{P}_{3}(S)$ » est écrit en interligne, au-dessus
de « l'ens.\ des »}
des arêtes $A \subset \mathfrak{P}_{2}(S)$ \add{est} défini comme l'ens.\ des
parties de card 2 contenues dans une face), par le choix \add{donc}
d'une orientation, \ill{}. Choisir cette orientation revient \ill{}
\ill{} $X = \mathbb{P}^{1}$, $S = \{0, \infty, \pm 1, \pm i\}$ et le choix
dans le cas standard \ill{} \uncertain{orientés} $\overleftrightarrow{\{0,1,i\}}$
et \ill{} une face est $\{0, 1, i\}$, et les faces
$\overrightarrow{\{0,1,i\}}$ sont \uncertain{opposées}. Donc le choix d'une
orientation revient au choix d'une racine carrée de $-1$.

Pour un sous-groupe $V \simeq D$ de $G_{X}$, on désigne par
$\operatorname{Oct}(V) \subset X$ l'octaèdre correspondant, par
$I(V) = F/\underline{a}_{I}$ l'ens.\ de ses paires de faces antipodiques. Alors
$\Gamma(V) = \operatorname{Nor}_{G}(V)$

\page{85}
qui opère sur $\operatorname{Oct}(V)$, puis sur $I(V)$, et on trouve
\[
  \Gamma(V) \xrightarrow{\ \sim\ } \mathfrak{S}_{I(V)} ,
\]
\note{sous $\Gamma(V)$, avec un signe d'égalité vertical, « \struck{\ill{}}
$\operatorname{Aut}^{+}(\operatorname{Oct}(V))$ »}
et est iso.\ \struck{\ill{}} définit une structure de $V$-torseur sur $I(V)$.
(Ceci résulte de la théorie combinatoire de l'octaèdre, ou dans \ill{} c'est
kif-kif.)

\textbf{Proposition} \struck{\ill{}} \uncertain{Soit} $X$ droite projective
\uncertain{réelle} sur un corps \ill{} dans le\supplied{quel} \ill{} soit
\uncertain{réalisable} par extraction de racines (il suffit qu'il
\struck{contienne} \ill{} \uncertain{racines} $\sqrt{\ }$ de car.\ $\neq 2$), et
$G = \operatorname{Aut}_{k} X$. Alors il y a corr.\ 1-1 entre
\begin{itemize}
  \item a) les sous-groupes $V$ de $G$ isom.\ à $D_{2}$
  \item b) — $\Gamma$ de $G$ isom.\ à $\mathfrak{S}_{4}$
    \struck{\ill{}}
  \item c) les parties $S \subset \mathfrak{P}_{6}(X(k))$ qui sont
    « octaédrales » i.e.\ \struck{telles que} \uncertain{isom.\ à}
    $(\mathbb{P}^{1}_{k}, \{0, \infty, \pm 1, \pm i\})$
\end{itemize}
\note{au b), le trait horizontal tient lieu de « les sous-groupes »}
par
\[
  \Gamma = \operatorname{Nor}_{G}(V)
\]
$V$ = unique sous-groupe \uncertain{invariant} de $\Gamma$ \uncertain{isomorphe}
à $D_{2}$
\note{« unique » est écrit au-dessus de « sous-groupe », et une fin de phrase
biffée suit « $D_{2}$ »}
$S$ = ens.\ des \uncertain{points fixes} de \uncertain{l'action} de $V$ sur
$X$ \struck{\ill{}}
(\emph{NB} $\Gamma$ a un ens.\ de points fixes plus grand qui comprend,
en plus de $S$, les « milieux d'arêtes » et « centres des faces » de
l'octaèdre),
\ill{} $S$, on a un antipodisme $\underline{a}_{S} = \underline{a}_{V}$ \ill{}
par la condition que pour $\sigma \in V^{\times}$, $\underline{a}_{S}$
\struck{est} \add{induit} la transposition \ill{} $X^{\sigma}$ (de cardinal 2).
\uncertain{On} définit sur $S$ une structure octaédrale, admettant $\Gamma$ comme
groupe des automorphismes orientés, et $\underline{a}$ comme antipodisme, via
l'antipodisme (les faces sont les parties $S' \in \mathfrak{P}_{3}(S)$ telles
que $S' \cup a(S') = S$).

\marginal{Cas : la catégorie formée des $(X, V)$ (ou $(X, \Gamma)$, ou
$(X, S)$) \ill{} \uncertain{équivalente} \ill{} des octaèdres \ill{}
$(\mathrm{Oc})$ \ill{} \ill{} Dans \ill{} \ill{}
$\mathbb{Q}(i)$ \ill{} $q : \{\pm i\} \simeq$ \ill{}}
\note{la note marginale est écrite en oblique, de bas en haut, sur toute la
moitié inférieure de la marge gauche ; redressée, elle reste en grande partie
illisible}

\page{86}
\emph{NB} On sait que la donnée d'un octaèdre \add{combin.,} \uncertain{non
orienté,} revient à celle d'un système $(I, \infty)$ \struck{\ill{}}
$E_{I_{3}} \times E_{I_{2}}$, ens.\ des \uncertain{code-graphes} et des
orientations. Ici $\omega = \{i, -i\}$, \struck{\ill{}} \struck{Mais ici on a
\uncertain{réduit} \ill{} $I$ \uncertain{figuré} \ill{} le donnée de $I$
\ill{}} qui \uncertain{s'était} \uncertain{ôté} \ill{}, mais \ill{} de
cardinal 6 \uncertain{sur lequel} $V$ opère. $S$, qui est un ens.\ de
\add{son} cardinal \ill{} \uncertain{définissent} une
\note{passage très raturé ; la ligne « Mais ici … figuré » est biffée
et la place de ce qui la remplace ne se décide pas}
La condition à mettre \struck{$*$} pour que \uncertain{ceci} définisse une
structure octaédrale (dont $S$ ser.\ l'ens.\ des sommets) \struck{\ill{}} est
que $\forall\, \sigma \in V^{*}$, \struck{\ill{}}
$\operatorname{card} S^{\sigma} = 2$, et
$S = \bigcup_{\sigma \in V^{*}} S^{\sigma}$ (réunion forcément disjointe) tel
\uncertain{condition} que $\sigma$ opère sur $(S \smallsetminus S^{\sigma})$
par transposition \add{\uncertain{dans les deux}} \struck{\ill{}} \ill{} des
composants $S^{\sigma'}$, $S^{\sigma''}$ (où $V^{*} = \{\sigma, \sigma',
\sigma''\}$). Donc la donnée de \struck{\ill{}} $(V, S, u : V \hookrightarrow
\mathfrak{S}_{S})$, ou $(S, V \subset \mathfrak{S}_{S})$, équivaut à celle de
$S$, avec une involution sans pts fixes \ill{} \ill{} \uncertain{ayant} les
$S^{\sigma}$ comme orbites).

\uncertain{Résumant} : la donnée d'un $V$ qui opère \add{fidt.} sur $X$,
considérons une orbite \add{$I$} de $V$ sur $X^{*} = X \smallsetminus S(V)$.
Donc $I$ est un torseur sous $V$. La donnée de $(X, I, V(I) \hookrightarrow
G_{X})$, $(X, I, I \hookrightarrow X)$ équivalent : celle de
$(X, I, V(I) \hookrightarrow G_{X}$, orbite de $V(I)$ dans $X)$, et on
\uncertain{se} \uncertain{garde} de confondre les $(V = V(I))$-torseurs $I$,
avec le $V$-torseur $I(V)$ \add{$= I(V, X)$} défini via l'octaèdre
$\operatorname{Oc}(V) = \operatorname{Oc}(V, X)$.

En termes de $I$ \add{$\{\pm i\}$}, on peut définir un octaèdre
\struck{orienté} \ill{} \struck{\ill{} bien \ill{}} canonique \struck{\ill{}
l'un des \ill{} « $\{\pm i\}$ » $I$, l'un de ses orientations \ill{}
\ill{}} \struck{bitransp.} \uncertain{à} condition que \ill{} \struck{soit
$V$}. Cet \struck{\ill{}} octaèdre, \struck{\ill{} pas choisi \ill{}
$i = \sqrt{-1}$}, peut \uncertain{s'écrire} comme ens.\ $(X, S)$ \ill{}
\struck{\ill{}} \uncertain{conséquemment} $(X(I), S(I))$. On a un \uncertain{iso}
\uncertain{canonique}
\note{les trois dernières lignes du feuillet sont chargées de ratures
superposées ; seules les formules et les mots isolés ci-dessus s'y lisent}

\page{87}
\[
  X(I)/V(I) \simeq X(J)
\]
où $J = J(I)$ comme dessus.

Ceci dit soit \add{une opération de $V$ sur $X$} \struck{$I \hookrightarrow X$
\ill{} donc \ill{} dans $X$ une} \struck{opération $V = V(I)$}
$V \hookrightarrow G_{X}$, d'où un $I(V)$ $V$-torseur \uncertain{des} \ill{}
\uncertain{paires} de $\operatorname{Oc}(V)$, et \ill{} \struck{$X \simeq$}
\[
  X \simeq X(I(V)) \qquad \text{iso.\ canonique}
\]
(et \uncertain{aussi} $X \simeq X(I)$). Mais \uncertain{pour} \struck{\ill{}} deux
torseurs $T, T'$ sous \ill{} $= V$, on a \ill{}
\[
  X(T \wedge T') = X(T) \wedge^{V} T'
\]
\marginal{\uncertain{explicitez}}
\note{à droite de la formule, derrière un trait vertical, un mot seul}
Soit donc $\underline{t} = \delta(I, \lambda)$ le $V$-torseur tel que
$I(V) \simeq I \wedge^{V} \underline{t}$, i.e.\
$\delta \simeq I(V) \wedge I^{-1}$, on aura donc
\[
  \boxed{X = X(I(V)) = X(I) \wedge^{V} \underline{t}}
\]
\struck{Donc} \add{donc} \struck{\ill{}}
\[
  \boxed{X/V \simeq X(I)/V \simeq \Sigma(J)}
\]
\note{la lettre de $\Sigma(J)$ est tracée comme un $X$ barré horizontalement ;
elle revient plus bas sous la forme « $\Sigma(J)^{*}$ »}

En résumé :

\emph{Cor.} \struck{\ill{}} \struck{La donnée, d'une} Pour
$I \in \operatorname{Ob} E_{I_{2}}$ \struck{\ill{}} donné, d'où $J$ \ill{}
et $V = V(J)$, la donnée d'un $X$ et d'un plongement
$V \hookrightarrow G_{X}$ équivaut à la donnée d'un torseur $\underline{t}$
sous $V$, via
\[
  \underline{t} \longmapsto X(I) \wedge^{V} \underline{t}
\]
On a donc
\[
  X/V \simeq \Sigma(J) \qquad \text{(iso.\ canonique)}
\]
\note{les trois paragraphes précédents, depuis « $V = V(J)$ », sont réunis
par un trait vertical dans la marge}

\struck{\ill{}} \add{Donc} la donnée de $(X, I \hookrightarrow X)$ équivaut à
celle de a) un pt $\xi$ de $\Sigma(J)^{*}$, b) un $V$-torseur $\underline{t}$,
\struck{et c) un $V$-isom.\ \ill{} $X = X(I) \wedge^{V} \underline{t}$, et}
\[
  X_{\xi} = X(I)_{\xi} \wedge \underline{t} \simeq I \qquad \text{i.e.} \qquad
  \underline{t} = I \wedge X(I)_{\xi}^{-1}
\]
\struck{$X_{\xi}$ = fibre de}

\page{88}
\textbf{Théorème} Soit $k$ un corps alg.\ clos de car.\ $\neq 2$, \struck{soit
$I$ un ens.\ de card 4, d'où $J$ de card.\ 3, et $V = V(J)$}
\struck{Considérons}
Considérons le groupoïde des $(X, I \hookrightarrow X)$, où $X$ est une droite
projective sur $k$, \add{$I$ de card.\ 4}. \struck{\ill{}} Ce groupoïde est
équivalent à celui des couples $(I, \xi \in \Sigma(J(I))^{*})$,
\uncertain{l'équivalence} \uncertain{étant} \ill{} : on pose
\[
  X = X(I) \wedge^{V(J(I))} t(\xi, I)
\]
\note{sous $V(J(I))$, une accolade et « $\simeq V(I)$ »}
où $t(\xi, I)$ est le $V(I)$-torseur $I \wedge X(I)_{\xi}^{-1}$
\[
  I \hookrightarrow X \text{ est la fibre de } X \to \xi, \text{ compte tenu
  que}
\]
\note{au-dessus de « la fibre », un $X$ en interligne, dont la place dans la
phrase n'est pas sûre}
\[
\begin{array}{c}
  X/V(I) \simeq X(I)/V(I) \simeq \Sigma(J) \\
  X_{\xi} \simeq X(I)_{\xi} \wedge t(\xi, I) .
\end{array}
\]
\note{l'énoncé, depuis « Théorème », est tenu par un trait vertical dans la
marge gauche}

\struck{\ill{}} Si $I$ est muni d'un pt $i_{0}$, on a une description
de $X$ comme $X(I) \wedge^{V(I)} X(I)_{\xi}^{-1}$, i.e.\ isom.\ à $X(I)$ une
fois choisi un \uncertain{élément} dans $X(I)$ au-dessus de $\xi$.

De façon plus générale, \struck{\ill{}} la donnée d'un isom.\
« admissible » $X \simeq X(I)$ équivaut à celle d'un $V$-isom.
\[
  I \simeq X(I)_{\xi} .
\]

\section{Courbes elliptiques réelles}
\note{titre souligné, en haut à droite du feuillet 90, de sa main}

\page{90}
\marginal{Déc 83}
\note{la date est écrite en oblique dans le coin supérieur gauche, séparée du
texte par un trait}

Soit $k$ un anneau (\uncertain{par exemple} commut.)
\struck{$M$ sur $k$-\ill{}}
Considérons la catégorie des $k$-Modules $M$ munis d'un automorphisme
$\sigma$ \struck{tel} tels que
\begin{itemize}
  \item a) $\sigma^{2} = \mathrm{id}$
  \item b) $2\,\mathrm{id}_{M}$ injectif (on dit \uncertain{alors} $M$
    \emph{2-régulier})
\end{itemize}
À un tel couple $(M, \sigma)$, associons les triples $(P, Q, m)$ de deux
\supplied{$k$-}modules pour lesquels $2\,\mathrm{id}_{P}$ et $2\,\mathrm{id}_{Q}$
sont injectifs (i.e.\ $P, Q$ 2-réguliers), et
$m \subset P/2P \oplus Q/2Q$ \add{un sous-$k/2k$-module} (avec
$m \cap P/2P = 0$, $m \cap Q/2Q = 0$, \uncertain{i.e.}\ $m$ est le graphe d'un
iso.\ entre deux sous-modules $p' \subset p = P/2P$, $q' \subset q = Q/2Q$, et
un iso.\ $p' \simeq q'$), de la façon suivante
\[
\begin{array}{l}
  P = M_{+} = \operatorname{Ker}(\mathrm{id}_{M} - \sigma)
    = \{x \in M \mid \sigma x = x\} \\
  Q = M_{-} = \operatorname{Ker}(\mathrm{id}_{M} + \sigma)
    = \{x \in M \mid \sigma x = -x\}
\end{array}
\]
\ill{} $M_{+} \cap M_{-} = 0$ dans $M$, car $M$ 2-régulier, \uncertain{donc}
\[
  2M \subset M_{+} \oplus M_{-} \subset M
  \quad \text{car } \forall x \in M, \text{ on a }
  2x = \underbrace{(x + \sigma x)}_{M_{+}} + \underbrace{(x - \sigma x)}_{M_{-}}
\]
\note{un point d'exclamation dans la marge, devant la formule ; à la suite
de « $\subset M$ », un terme biffé : \struck{$\subset \tfrac{1}{2}(M_{+} \oplus M_{-})$}}
\uncertain{soit}
\[
  \underset{\substack{\Vert \\ 2M}}{M'} \subset M_{+} \oplus M_{-} = P \oplus Q
\]
on a donc
\struck{\ill{}}
\[
  2P \oplus 2Q \subset M' \subset P \oplus Q
\]
donc $M'$ est l'image d'un sous-module $m$ de
\[
  (P \oplus Q)/(2P \oplus 2Q) \simeq \underbrace{P/2P}_{p} \oplus
  \underbrace{Q/2Q}_{q} ,
\]
et on \uncertain{vérifie}
\[
\begin{cases}
  m \cap p = 0 & \text{i.e.\ }
    \underbrace{M' \cap \bigl(P + (2P + 2Q)\bigr)}_{M' \cap P + (2P + 2Q)}
    = (2P + 2Q) \\
  m \cap q = 0 & \text{i.e.\ } M' \cap \bigl(Q + (2P + 2Q)\bigr) = 2Q
\end{cases}
\]
car $M' \cap P = 2P$ puisque $2x \in P$ i.e.\ $\sigma 2x = 2x$ implique
$x \in P$ i.e.\ $\sigma x = x$, \uncertain{et} $M' \cap Q = 2Q$.

\uncertain{D'ailleurs} \struck{\ill{}} $(M, \sigma)$ se reconstitue à partir
de $(P, Q, m)$, car $M' \subset P \oplus Q$ se reconstitue, et
\struck{\ill{}} $\sigma_{M'}$ est induit par
$\mathrm{id}_{P} \oplus (-\mathrm{id}_{Q})$, et
$2\,\mathrm{id}_{M} : M \to M'$ est un iso.\ de $(M, \sigma_{M})$ sur
$(M', \sigma_{M'})$.

\emph{Prop.} Le foncteur précédent $(M, \sigma) \longmapsto (P, Q, m)$ est une
équivalence de catégories.

\emph{Cor.} Supposons $k$ \struck{\ill{}} \add{connexe,} 2-régulier, prenons
\struck{\ill{}} $M$ \add{$(M, \sigma)$} \uncertain{projectif} de
\uncertain{rang} 2 \ill{}, avec $\sigma \neq \mathrm{id}_{M}, -\mathrm{id}_{M}$,
alors la cat.\ des \struck{\ill{}} \uncertain{des} \uncertain{faisceaux}
\ill{} \uncertain{est} équivalente à celle des paires $(P, Q)$ avec $P, Q$
inv., et $p \subset P/2P$, $q \subset Q/2Q$ et $p \simeq q$ \ill{}
\note{la dernière ligne est coupée au bord droit du feuillet ; l'interligne au-dessus de « prenons » porte plusieurs mots biffés}

\page{91}
\note{calculs de brouillon pour le feuillet précédent, écrits sur un
listing d'ordinateur daté « 19 AUG 82 » ; ils sont transcrits dans l'ordre de
la page, du haut à droite vers le bas}

\struck{$M'$}
\[
  P \oplus Q \supset M' \supset 2P \oplus 2Q
\]
\struck{\ill{} $\exists\, x, y \in M'$, $(1 + \sigma)(x, y) =$ \ill{}}
\[
\begin{array}{l}
  \sigma(x, y) = x - y \\
  (1 + \sigma)(x, y) = 2x \\
  (1 - \sigma)(x, y) = 2y
\end{array}
\]
\note{la troisième ligne est barrée d'un trait oblique ; le « $x - y$ » de la
première est tel qu'écrit}
$P =$ \struck{$x \in (x, y) \in M'$}
\[
\begin{cases}
  M' \cap P = 2P \\
  M' \cap Q = 2Q
\end{cases}
\qquad \text{i.e.\ } M' \cap (P + 2Q) = 2P + 2Q
\ \ \bigl(= (M' \cap P) + 2Q\bigr)
\]
$(x, y) \mapsto x + 2y \in M'$
\[
  m \subset \underbrace{P/2P}_{p} \times \underbrace{Q/2Q}_{q}
  \qquad \text{sous-module de } k/2k\text{-modules}
\]
\[
\begin{array}{ll}
  m \cap p = 0 & \qquad m \to p \text{ injectif} \\
  m \cap q = 0 & \qquad m \to q \text{ injectif}
\end{array}
\]

$x \mapsto 2x : M \to M$ injectif \quad
\struck{$k'$} $k \hookrightarrow k' = k[\tfrac{1}{2}]$

$M$ plat sur $k$ \quad
$M \subset M \otimes k' = M' = M'_{+} \oplus M'_{-}$ \struck{$= M_{+} k'$}
\[
\begin{array}{l}
  M_{+} = \{x \in M \mid \sigma x = x\} = \operatorname{Ker}(1 - \sigma) \\
  M_{-} = \{x \in M \mid \sigma x = -x\} = \operatorname{Ker}(1 + \sigma) \\
  M_{+} \cap M_{-} = 0 \\
  M/M_{+} + M_{-} \ \text{annulé par } 2 \\
  2x = \underbrace{(x + \sigma x)}_{\in M_{+}} + \underbrace{(x - \sigma x)}_{\in M_{-}}
\end{array}
\]
dans \uncertain{$x \mapsto 2x$} $M \to M$ injectif
\[
\begin{array}{l}
  M_{+} \oplus M_{-} \subset M \subset \tfrac{1}{2}(M_{+} \oplus M_{-}) \\
  2M_{+} \oplus 2M_{-} \subset \underbrace{2M}_{M'} \subset M_{+} \oplus M_{-}
\end{array}
\]
\struck{$P \oplus Q$} \ $P, Q$ \qquad $M$
\[
  m \in P/2P \oplus Q/2Q
\]
\[
  (1 + \sigma) M' = 2M_{+}, \qquad (1 - \sigma) M' = 2M_{-}
\]
\note{un long trait oblique traverse la moitié inférieure de ces calculs, de
« $M_{-} = \ldots$ » jusqu'à « $2M$ » : la page est ainsi, en partie,
cancellée}

\page{92}
\emph{Courbes elliptiques réelles} \struck{Elles sont} La catégorie de ces
courbes (avec morphismes \uncertain{pas} \uncertain{nécessairement} des iso)
\uncertain{équivaut} à celle des systèmes \struck{\ill{}}
$(\Pi, \sigma, c)$, où $\Pi$ est un $\mathbb{Z}$-module libre \add{de rang 2},
\struck{\ill{}} $c$ une structure \struck{conf.} conforme orientée sur
$\Pi_{\mathbb{R}} = \Pi \otimes_{\mathbb{Z}} \mathbb{R}$, et $\sigma$ un
anti-automorphisme de \uncertain{la} \uncertain{struct.} $(\Pi, c)$, qui peut
s'interpréter comme une \add{\uncertain{symétrie} de $\Pi \otimes \mathbb{R}$
i.e.\ de déterminant $-1$,} involution de $\Pi$, renversant l'orientation
(donc $\neq \pm \mathrm{id}_{\Pi}$), et conservant la structure conforme. Pour
le couple $(\Pi, \sigma)$ il y a seulement deux possibilités, à
\uncertain{iso} près (elles sont d'ailleurs \uncertain{isogènes} …). Je
\uncertain{note}, pour chacune d'elles, \uncertain{à} trouver les structures
conformes qui sont invariantes par $\sigma$. On peut les normaliser par
\struck{\ill{}}
\[
  Q(e_{+}) = 1
\]
\note{après $Q(e_{+})$, une parenthèse biffée ; en marge droite, « \emph{NB}
\ill{} une base de \struck{\ill{}} $\Pi_{+}$, $\Pi_{-}$ », avec $\Pi_{+}$ et
$\Pi_{-}$ l'un sous l'autre et un trait qui les relie}
donc on sait les structures de la forme
\[
  Q(x e_{+} + y e_{-}) = x^{2} + a y^{2}
\]
où $a \in \mathbb{R}^{*+}$. Il y a donc deux familles de courbes elliptiques
réelles, paramétrées chacune par $\mathbb{R}^{*+}$, on peut les désigner par
$E_{a}$, $E'_{a}$ $(a \in \mathbb{R}^{*+})$. \struck{Considérons}

Considérons les automorphismes \struck{\ill{}} \add{d'une courbe elliptique}
définis par $(\Pi, \sigma, c)$, ils correspondent \struck{\ill{}} \add{aux}
\uncertain{automorphismes} \ill{} \uncertain{commutant} à $\sigma$ et
\uncertain{respectant} la \ill{} de $\Pi$ qui \add{\uncertain{sont}
\uncertain{orientés}} \struck{\ill{}} \ill{} = \uncertain{celles} de la
structure conforme (\uncertain{compatible} à l'orientation) \struck{\ill{}}
\ill{} $Q$ transforme $\Pi^{\pm}$ en $\Pi^{\pm}$ \struck{\ill{}}, \ill{}
\ill{} $\pm e_{0}$, \ill{} \ill{} \struck{\ill{}} \ill{} la symétrie
\ill{}) \ill{} \ill{}. Mais \ill{} \ill{} \ill{} de
$\mathbb{R} e_{0} + \mathbb{R} e_{1}$ \uncertain{qui respectent} $Q$, qui
\uncertain{fixent} $e_{0}$ et qui \ill{} \uncertain{respectent}
l'orientation ! Donc il \uncertain{n'y a} que l'identité, et $\sigma$,
\struck{\ill{}} \ill{} que \uncertain{l'automorphisme} universel,
\uncertain{correspondant} à $-1$.
\note{la seconde moitié du feuillet est une écriture rapide, très raturée ;
les formules et quelques mots s'y lisent, la prose presque pas}

\page{93}
Pour une courbe elliptique \uncertain{complexe} donnée, définie par un réseau
\add{$\Pi$} dans un espace vectoriel complexe $V$ de dim 1, les structures
réelles qu'on peut mettre dessus correspondent aux \struck{\ill{}}
\add{droites réelles} \ill{} de $V$, telles que la \struck{symétrie}
\add{réflexion} \uncertain{orthogonale} \ill{} : $\mathbb{D}$ conserve $\Pi$.
Peut-il y en avoir plusieurs ? \uncertain{Si oui}, les structures réelles
correspondantes sont-elles isom.\ (i.e.\ conjuguées de l'une d'elles par
\ill{} automorphisme complexe \ill{}) ? \struck{En tous cas, \ill{}
$\sigma, \sigma'$ \ill{}}

\struck{Cas de} \struck{\ill{}} Dans tous les cas, soit $G$ le groupe des
\uncertain{automorphismes} \add{$X$ défini par $(\Pi, c)$} de la courbe
elliptique \uncertain{conformes}. Pour que \struck{\ill{}} $X$ admette une
structure réelle, il faut \struck{\ill{}} \uncertain{et il suffit} que
$G \xrightarrow{\det} \pm 1$ soit surjectif, \ill{} \ill{} \ill{} : on a
\[
  1 \longrightarrow \underset{\substack{\Vert \\ \operatorname{Aut}(X)}}{G^{+}}
  \longrightarrow G \longrightarrow \pm 1 \longrightarrow 1
\]
Mais un élément $\sigma$ de $G$ de déterminant $-1$ est une transformation
orth.\ de déterminant $-1$, et \uncertain{une} réflexion, donc c'est une
structure réelle sur la courbe elliptique ; les courbes elliptiques réelles
\uncertain{sur} la \ill{} \uncertain{structures} réelles \ill{} \ill{} qui
\uncertain{admettent} une \ill{} \ill{} \uncertain{celles} dont l'invariant
\ill{} \ill{} \uncertain{sous} la conjugaison, i.e.\ \struck{\ill{}} $G^{+}$ est
cyclique d'ordre \struck{\ill{}} $n = 2, 4$ \uncertain{ou} $6$ et
\add{\uncertain{réel}}. L'ens.\ \struck{des \ill{} structures} \struck{\ill{}}
\ill{} est un \uncertain{torseur} sous $G^{+}$, \struck{\ill{} ($X$ fini)}
et $G$ \ill{} \struck{\ill{}} \uncertain{isomorphe} au groupe diédral
$D_{n}$. $G^{+}$ \ill{}
\note{le bas du feuillet est chargé de ratures en deux couches ; « $n = 2, 4,
\ldots$ » est ajouté au-dessus d'une ligne biffée, « \ill{} diédral $D_{n}$ »
termine la page}

\page{94}
structures réelles s'identifie ! $G^{-}$, qui est un \emph{torseur} sous
$G^{+}$ par translations à gauche. Par contre, l'opération de $G^{+}$ sur
$G^{-}$ par automorphismes intérieurs est \ill{} \ill{}, le centre
$\zeta = \{\pm 1\}$ de $G$ \uncertain{opérant} trivialement, et il y a deux
orbites sous $G^{+}$ (ou sous $G^{+}/\zeta$) \uncertain{si} $n = 2$
\uncertain{ou} \ill{} (\uncertain{pour} $n = 4$ \uncertain{ou} $6$) les
symétries par rapport \struck{\ill{}} \add{respectivement} aux sommets et aux
côtés des \struck{\ill{}} polygones (carrés ou hexagones) fondamentaux
\note{dans la marge gauche, à hauteur de cette ligne, un petit carré dessiné}

\emph{Donc} on \uncertain{trouve}

\emph{Proposition}. Soit une courbe elliptique \struck{réelle} complexe $X$.
Les structures réelles \struck{\ill{}} \add{sur $X$} \struck{\ill{}}
\ill{} de la courbe \struck{\ill{}} \ill{} si et seulement \ill{}
\struck{\ill{}}
\[
  1 \longrightarrow G^{+} \longrightarrow G \longrightarrow \pm 1
  \longrightarrow 1
\]
\ill{} et carrés (i.e.\ \uncertain{droites}). Dans ce cas, $G$ est
\struck{un groupe} \add{donc un} groupe diédral, dont $G^{+}$ (d'ordre
$n = 2, 4, \ldots, 6$) est le groupe des rotations. L'ens.\ des structures
réelles se \uncertain{partage} en deux orbites sous $G^{+}$ (\uncertain{par}
\ill{} \ill{}), et les structures réelles correspondant à deux orbites
distinctes sont (par définition !) non isomorphes. \struck{\ill{} \ill{}
correspondant aux structures $E_{a}$, $E'_{a}$ \ill{}}

Je présume que l'une est formée de structures $E_{a}$
($a \in \mathbb{R}^{*+}$ fixé), l'autre est formée des structures $E'_{a'}$
(avec $a' = a$), et inversement, i.e.\ qu'on peut trouver un
$a \in \mathbb{R}^{*+}$ unique, tel que les deux orbites sont l'une formée des
structures isom.\ à $E_{a}$, l'autre des structures isom.\ à $E'_{a}$. De plus
il faudrait \ill{}

\page{95}
\note{feuillet de dessins, sans un mot ni une formule : neuf sphères tracées
à main levée au feutre (noir, brun, vert, orange, bleu), chacune découpée par
un réseau d'arcs de grands cercles en régions dont certaines sont hachurées
ou coloriées ; la plus grande, en orange, en haut à droite, est coupée par
quatre cordes en secteurs, l'un hachuré}

\page{96}
explicitons le lien \uncertain{du} paramètre $a$ avec le paramètre $\tau$, ou
mieux, avec l'invariant modulaire $j \in \mathbb{R}^{\pm}$
(\uncertain{s'agit-il} \struck{i.e.\ $a = \exp j$ ou $a = \exp - j$}
\struck{\ill{} dans la famille $E_{a}$ ou $E'_{a}$ ?)}).
\marginal{$j = \tfrac{1}{2} \log a$ ?}
\note{dans la marge gauche, sous cette note, un hexagone dessiné avec son axe
vertical et une diagonale horizontale}

\struck{Considérons} Notons que si $\sigma$ est dans $G^{-}$, \uncertain{alors}
(\uncertain{réflexion} par \ill{} \ill{} \uncertain{hexagonal}, \ill{})
suivant les \uncertain{cas} $-\sigma$ est conjugué, ou non conjugué, à
$\sigma$. Dans le cas hexagonal il est non conjugué ; dans le cas carré, il
est conjugué. Dans le cas \uncertain{non harmonique}, $n = 2$, il est
\uncertain{bien sûr} non conjugué (cas \uncertain{d'un} \uncertain{groupe}
commutatif) -- dans le cas \struck{hexagonal} \add{carré} et le seul \ill{}
\ill{} \struck{\ill{}} soit conjugué à $\sigma$. \struck{Dans le cas \ill{}}
\struck{dans le cas carré précisément, on trouve}

\emph{Distinguons les cas}

1) \emph{$n = 2$} Cas \uncertain{non harmonique}. Le groupe des automorphismes
\add{anti} est $\pm 1$, $\pm \sigma$, où $\sigma$ est une \uncertain{anti-rotation}.
On voit donc que si $\sigma$ est du type \ill{} \ill{} droite d'un
\ill{}, il en est de même de $-\sigma$ (\ill{} $P$, $Q$ \ill{}
$\mathrm{id}_{P}$ et $-\mathrm{id}_{Q}$), donc on est dans le cas $E_{a}$,
\ill{} \ill{}) donc on est dans $E_{a}$ et en $E_{a'}$, correspondant aux
deux \ill{} de \ill{} $a a' = 1$. \struck{\ill{}} \uncertain{Si} on
\uncertain{prend} $\Pi$, $\sigma'$ de $G^{-}$, \ill{} \ill{}
\uncertain{un générateur} $u = e_{+}$ de $\Pi_{+}$, \ill{} $v = e_{-}$ de
$\Pi_{-}$, \ill{}
\[
  a = a(\sigma) = \frac{Q(e_{-})}{Q(e_{+})}
\]
et \uncertain{remplaçant} \ill{} la \ill{} de $\sigma$ \uncertain{par} $-\sigma$
\ill{} \ill{}, on trouve bien
\[
  a(-\sigma) = 1/a(\sigma)
\]
\marginal{\uncertain{Ceci} vaut \uncertain{aussi} bien dans le cas des
$E_{a}$ que des $E'_{a}$}
(valable dans tous les cas), donc
\[
  a(-\sigma) = a(\sigma) \iff a(\sigma) = 1 \iff \text{cas \emph{carré}} .
\]

\page{97}
2) \emph{Cas carré} $n = 4$. On voit que l'un \uncertain{a} un $\sigma$ du type
$E_{1}$, et $-\sigma$ est du \uncertain{même} type. Cela \uncertain{correspond}
aux symétries p.r.\ aux \add{diagonales} \struck{\ill{}} des carrés de
\uncertain{départ}. La symétrie par rapport à une \uncertain{arête},
\uncertain{changeant} $e_{0}$ et $e_{1}$, les \uncertain{invariants} et
anti-invariants \uncertain{sont} engendrés par $e_{0} + e_{1}$ et
$e_{0} - e_{1}$, \uncertain{qui engendrent} \ill{} \ill{} \ill{} dans le cas
$E'_{a'}$. \uncertain{Comme}
\[
  a' = \frac{Q(e_{0} - e_{1})}{Q(e_{0} + e_{1})} = 1
\]
\note{devant le « 1 », une lettre biffée : \struck{A}}
on \uncertain{voit} \ill{} \ill{}. Donc on trouve les types $E_{1}$
(« \uncertain{sommets} ») et $E'_{1}$ (« côtés »).
\note{dans la marge gauche, le carré fondamental dessiné sur ses axes, avec
les vecteurs $e_{0}$ et $e_{1}$ issus du centre vers deux sommets}

3) \emph{Cas hexagonal} $n = 6$.

\struck{\ill{}} \emph{Réflexions-sommets}. L'\uncertain{ensemble} des
\uncertain{invariants} \uncertain{par} la réflexion p.r.\ à $\mathbb{R} u$ sont
engendrés par $u$, les anti-invariants par $2v - u$, c'est du type $E'_{a}$,
\uncertain{avec}
\[
  a = \frac{Q(2v - u)}{Q(u)} = \sqrt{3}
\]
\note{dans la marge gauche, l'hexagone dessiné avec ses six sommets nommés
$u$, $v$, $v - u$, $-u$, $-v$, $u - v$, et le sommet $2v - u$ du triangle
extérieur qui s'appuie sur le côté $v - u$, $v$ ; des flèches vont du centre
aux sommets. La valeur $\sqrt{3}$ est telle qu'écrite}

\emph{Réflexions-faces} \uncertain{Ce sont} les $-\sigma$, on trouve donc
$E'_{1/a} = E'_{1/\sqrt{3}}$

\emph{En résumé} Dans tous les cas, sauf le cas carré, \struck{\ill{}}
\add{l'un des}
\struck{les types de} \struck{les} structures réelles sur $X$ est soit
\struck{du type} $\{E_{a}, E_{1/a}\}$, \add{avec $a > 1$,} \struck{Cas} soit
$\{E'_{a}, E'_{1/a}\}$ (le cas hexagonal étant dans le deuxième cas \add{avec
$a = \sqrt{3}$}), à l'exclusion bien sûr l'un de l'autre.
\emph{Dans} le cas carré, l'ens.\ des structures réelles est
$\{E_{1}, E'_{1}\}$ i.e.\ il y a à la fois des structures $E$ et $E'$.
\note{le résumé est tenu par un trait vertical dans la marge gauche}

\page{98}
\emph{Résumé}

Cas $E_{a}$ : base $u = e_{0}$, $v = e_{1}$, forme
\struck{$Q(e_{0}) = Q(e_{1})$} $Q(x e_{0} + y e_{1}) = x^{2} + a y^{2}$.

Si on \uncertain{veut} \struck{\ill{}} \add{\uncertain{normaliser}} forme $Q$
sous forme standard, on prend la base
$\underbrace{e_{0}}_{1}, \underbrace{\tfrac{1}{\sqrt{a}}\, e_{1}}_{i}$, on
trouve donc la base $1, \sqrt{a}\, i$ d'où
\[
  \boxed{\tau = \sqrt{a}\, i}
\]

Cas $E'_{a}$ : base \struck{$e_{0}$, $v =$ \ill{} $u$, \ill{}} \ill{} \ill{}
\ill{} \ill{} la base $e_{0} = u$, $e_{1} = 2v - u$, où
$Q(x e_{0} + y e_{1}) = x^{2} + a y^{2}$, $v = \tfrac{1}{2}(e_{0} + e_{1})$ ;
on \struck{\ill{}} \uncertain{change} \uncertain{encore} de base.
\[
  \underbrace{e_{0}}_{1}, \ \underbrace{\tfrac{1}{\sqrt{a}}\, e_{1}}_{i}
  \qquad \text{base orthonormale}
\]
et on a $u \mapsto 1$, $v \mapsto \tfrac{1}{2}(1 + \sqrt{a}\, i)$, donc
\[
  \boxed{\tau = \tfrac{1}{2}(1 + \sqrt{a}\, i)} = \frac{1}{2} + \frac{\sqrt{a}}{2}\, i
\]
\note{un trait ondulé traverse ici toute la largeur du feuillet}

Calcul des groupes des pts réels d'une courbe $E_{a}$, $E'_{a}$.

Pts réels sont les $x \bmod \Pi$ $(x \in \Pi_{\mathbb{R}})$ tels que
\[
  \sigma x \equiv x \quad (\Pi)
\]

\emph{Distinguons les deux cas}.

1) \emph{$E_{a}$} \struck{$\sigma(e_{0}, e_{1}) =$} $\sigma(x, y) = x, -y$
donc
\[
  \sigma(x, y) - (x, y) = (0, -2y)
\]
\note{le signe de « $-2y$ » est un trait court devant le 2, dont la lecture
n'est pas sûre}
donc la condition est que $2y \in \mathbb{Z}$. Donc $y = 0$ \struck{ou}
\add{ou $\tfrac{1}{2}$}, les pts de
\struck{$0$,} $\tfrac{1}{2}\, \mathbb{R} e_{0}$ et de
\struck{\ill{}} $(0, \tfrac{1}{2}) + \mathbb{R} e_{0}$, donc \emph{deux
composantes connexes} pour $X(\mathbb{R})$.

\page{99}
\emph{Cor.} Supposons $k$ principal \add{2-régulier}, $k_{0} = k/2k$ un corps.
Alors les $(M, \sigma)$ \add{\uncertain{à} \uncertain{isom.\ près}}
\uncertain{projectifs} équivalent à celles des paires
$(P, Q,$ \struck{\ill{}}$)$ avec $P, Q$ libres de rang 1, et
\struck{$p \subset$} $m \subset \underbrace{P/2P}_{p} \oplus
\underbrace{Q/2Q}_{q}$ sous-esp.\ vectoriel sur $k_{0}$, \struck{\ill{}}
qui est soit nul, soit une droite $\neq p, q$ (\uncertain{toutes}
\uncertain{conditions}).

\emph{Cor.} Cela donne une description des classes de conjugaison
d'éléments de $\mathrm{Gl}(2, k)$ \struck{$*$} qui sont d'ordre 2 exactement,
\add{\uncertain{i.e.}} et distincts de $\{\pm 1\}$ : il y \struck{\ill{}} a
$\begin{pmatrix} 1 & 0 \\ 0 & -1 \end{pmatrix}$ (correspondant au cas
$m = 0$), puis \struck{\ill{} \uncertain{familles} indexées par \ill{} \ill{}
\ill{} $\mathbb{P}^{1}(k_{0})$ distincts de $0, \infty$ (\ill{} \ill{} que
$k^{*} \to k_{0}^{*}$ est surjectif) \ill{} \ill{} la classification \ill{}
plus \ill{})} \add{$k^{*} \to k_{0}^{*}$ \uncertain{surj.}} \uncertain{une}
\uncertain{fois} \ill{} base de $P/2P$ \ill{} $Q/2Q$ \uncertain{provenant}
\uncertain{d'}une base de $P$, $Q$.
\note{deux lignes de ce paragraphe sont biffées ensemble et reprises dans la
marge droite par « $k^{*} \to k_{0}^{*}$ surj. »}

\uncertain{Elle} correspond à
\[
\begin{array}{l}
  P = k, \quad Q = k, \quad m = \text{diagonale de } k_{0} \times k_{0}
    \text{ i.e.} \\
  M' = (2k \oplus 2k) + \text{diag.\ de } k \oplus k
    = \{(x, y) \mid x - y \in 2k\}
\end{array}
\]
donc \uncertain{les} \uncertain{pour} prendre base de $M'$
\[
  u_{0} = (2, 0), \quad v_{1} = (1, 1), \quad \text{d'où} \quad
  -u_{0} + 2 v_{1} = (0, 2) = \operatorname{sym}(u_{0})
\]
\note{les lettres de $u_{0}$ et $v_{1}$ sont surchargées à chacune de leurs
occurrences sur cette ligne et la suivante ; la lecture $u$, $v$ est celle
que confirme la figure en marge}
\marginal{$u = 2e_{0}$, $v = e_{0} + e_{1}$}
On trouve
\[
  \sigma u_{0} = u_{0}, \quad \sigma v_{1} = (1, -1) = u_{0} - \operatorname{sym} v_{1}
  = u_{0} - v_{1}
\]
i.e.\ la matrice de $\sigma$ p.r.\ à $(u_{0}, v_{1})$ est
\[
  \begin{pmatrix} 1 & 1 \\ 0 & -1 \end{pmatrix}
\]
Dans le cas où $k = \mathbb{Z}$, c'est la matrice de la réflexion p.r.\ à la
diagonale $u, e_{0}$ de l'hexagone régulier, \uncertain{hexagone} \ill{}, le
\ill{} de l'hexagone étant réalisé à un \uncertain{rapport} de
\uncertain{similitude} \uncertain{près} \ill{} \uncertain{comme} base deux
sommets consécutifs $u$, $v$ (\uncertain{Mais} le réseau \uncertain{est}
engendré par \ill{} un des sommets).
\note{dans la marge gauche, l'hexagone dessiné, sommets nommés
$v - u$, $v$, $-u$, $e_{0} = u_{+}$, $-v$, $\sigma v = u - v$, avec en haut
« $2v - u = u_{-}$ » ; flèches du centre vers les sommets}

\emph{NB} Soit $u_{+}$ une base de $\Pi_{+}$, $u_{-}$ une base de $\Pi_{-}$,
alors \struck{dans} on peut représenter $\sigma_{\mathbb{R}}$, pour la base
$(u_{+}, u_{-})$ de $\Pi \otimes \mathbb{R}$, par la matrice
$\begin{pmatrix} 1 & 0 \\ 0 & -1 \end{pmatrix}$, le réseau $\Pi$ étant
identifié à $\mathbb{Z} u_{-} + \mathbb{Z} u_{+}$ \struck{\ill{}} dans le
premier cas et dans le $2^{\text{e}}$ \uncertain{cas} à l'ens.\ des
$x u_{+} + y u_{-}$ tels que $x - y \in \mathbb{Z}$, $x, y \in \tfrac{1}{2}\mathbb{Z}$
\note{la dernière ligne est coupée au bord droit ; au pied du feuillet, écrit
serré sous la dernière ligne : « $\sigma^{2}(e_{1}) = \sigma e_{0} - \sigma
e_{1} = e_{0} - (e_{0} - e_{1}) = e_{1}$ »}

\page{100}
2) \emph{$E'_{a}$} Donc $\Pi$ engendré par $u, v$ avec $\sigma u = u$,
$\sigma v = u - v$
\note{dans la première ligne de calcul, le second $\sigma$ est surchargé ou
biffé}
\[
\begin{array}{l}
  \sigma(x, y) = \sigma(xu + yv) = xu + y(u - v) = (x + y) u - y v . \\
  \sigma(x, y) - (x, y) = yu - 2yv = y(u - 2v),
\end{array}
\]
donc la condition est $y \in \mathbb{Z}$, donc on trouve \struck{\ill{}} les
points de $\mathbb{R} u_{+}$ mod $\Pi$, donc dans ce cas $X(\mathbb{R})$ est
connexe.

\marginal{9.2.}
\note{« 9.2. » est écrit dans la marge gauche, à hauteur du « NB » ; c'est
peut-être une date, sans année}
\emph{NB} Pour bien faire, il faudrait aussi classifier les courbes elliptiques
réelles inhomogènes, i.e.\ les torseurs $X$ sur $\operatorname{Spec}
\mathbb{R}$, sous les \add{schémas en} groupes elliptiques $E_{a}$, $E'_{a}$.
Les torseurs non triviaux fournissent les surfaces conformes \struck{sans}
bord non orientables de genre 1.
\note{« sans » paraît barré d'un trait horizontal, et aucun mot n'est écrit
à sa place ni en interligne ; la lecture « à » d'une première passe ne se
trouve pas sur la page} La suite exacte de
$\gamma$-modules (où $\gamma = \operatorname{Gal}(\mathbb{C}/\mathbb{R})
\simeq \{\pm 1\}$)
\[
  0 \longrightarrow \Pi \longrightarrow V \longrightarrow
  E(\mathbb{C}) \longrightarrow 0
\]
fournit
\[
  H^{1}(\gamma, E(\mathbb{C})) \simeq H^{2}(\gamma, \Pi) \simeq
  \Pi^{\gamma}/(1 + \sigma)\Pi = \Pi_{+}/(1 + \sigma)\Pi
\]

1) \emph{Cas de $E_{a}$}. On a $\Pi \simeq \Pi_{+} \oplus \Pi_{-}$, d'où
$\Pi^{\gamma} = \Pi_{+}$
\[
  H^{1}(\mathbb{R}, E_{\mathbb{R}}) \simeq H^{1}(\gamma, E(\mathbb{C}))
  \simeq H^{2}(\gamma, \Pi) \simeq
  \underbrace{H^{2}(\gamma, \Pi_{+})}_{\mathbb{Z}/2\mathbb{Z}} \oplus
  \underbrace{H^{2}(\gamma, \Pi_{-})}_{0}
\]
\note{sous chacun des deux termes, un signe d'égalité vertical porte la
valeur}
Il y a exactement une classe non nulle dans $H^{1}$, qui est d'ordre 2

2) \emph{Cas de $E'_{a}$} On a $H^{1}(\mathbb{R}, E_{\mathbb{R}}) \simeq 0$.

Donc les seules courbes elliptiques vraiment inhomogènes sur $\mathbb{R}$ sont
les $E^{!}_{a}$ = \add{le} torseur \struck{\ill{}} non trivial (d'ordre 2) sous
$E_{a}$, dépendant du paramètre $a \in \mathbb{R}^{*+}$.
\note{le « ! » en exposant de $E^{!}_{a}$ est tel qu'écrit}

\end{document}
